\documentclass{article}
\usepackage{amsmath}

\title{Asignación 2 -  Tema I.1, Concepto de Curva}
\author{Elaborado por David Duarte}
\date{\today}

\begin{document}
	\maketitle
	
	\section*{Ejercicios}
	
	\begin{enumerate}
		\item Encuentra la longitud de arco de la curva paramétrica $\mathbf{r}(t) = (t, t^2)$ en el intervalo $0 \leq t \leq 1$.
		\item Calcula la curvatura de la curva $\mathbf{r}(t) = (\cos t, \sin t, t)$ en el punto donde $t = \pi/4$.
		\item Demuestra que la curva $\mathbf{r}(t) = (t, t^3, t^2)$ es una curva plana.
		\item Calcula la longitud de arco de la espiral logarítmica $\mathbf{r}(t) = (ae^{bt}\cos t, ae^{bt}\sin t)$ en el intervalo $0 \leq t \leq 2\pi$.
		
		\item Calcula la curvatura de la hélice circular recta $\mathbf{r}(t) = (a\cos t, a\sin t, bt)$.
		\item Encuentra el vector ortonormal de la curva $\mathbf{r}(t) = (t, t^2, t^3)$.
		\item Demuestra que la curva $\mathbf{r}(t) = (a\cos t, a\sin t, bt)$ es una curva plana.
		\item Calcula la longitud de arco de la curva $\mathbf{r}(t) = (e^t, e^{-t}, 2t)$ en el intervalo $0 \leq t \leq 1$.
	\end{enumerate}
	Elabore un código para cada ejercicio, con su respectivo gráfico, en el lenguaje de su preferencia.
\end{document}
